% section: おまけ：定数数列

\section{おまけ：定数数列}
数列の種類には定数数列というものがある.
時折,漸化式を解く中で定数数列が生まれることがある.
先程の解説から引用しよう.
\[
  \frac{a_{n}}{n(n+1)}=\frac{a_{n-1}}{(n-1)(n)}=\frac{a_{n-2}}{(n-2)(n-1)}=\cdots=\frac{a_{1}}{1(1+1)}
\]
これはまさに定数数列の良い例である.
少し式が複雑なので,
\[
  b_n=\frac{a_{n}}{n(n+1)}
\]
とおくと式は
\[
  b_n=b_{n-1}=b_{n-2}=\cdots=b_1
\]
となる.
ずっと同じ値を取るのだ.
当たり前過ぎて何を言っているんだ,と思うかもしれないが急に出てきてパニックにならないようにここで念のため触れておく.
